\section{Historical Notes}

\subsection{The Number $e$}

\begin{figure}
  \begin{center}
  \includegraphics[width=10cm]{napier_tables.jpg}
  \end{center}
  \caption{A page from the tables calculated by John Napier
  with the values of the $\sin$ function (Sinus) and its logarithm
  (Logarithmi).
  This page contains the values 
    for angles from 9 degrees (Gr.) to 9 and a half degrees (sexagesimal degrees
  are divided into 60 minutes) and
  their complements from 80 and a half degrees to 81 degrees.
  The code written on page~\pageref{code:napier} allows
  the calculation of the same values in a few hundredths of a second,
  also highlighting some errors in the last decimal digits.
  Napier took 20 years to complete the compilation 
  of the tables from 0 to 90 degrees.
    For example, if we wanted to calculate the value of the tangent of
  $\alpha=9$ degrees,
  we could observe that the logarithm of $\tg \alpha$ 
  is the difference between the logarithm of $\sin \alpha$ and the 
  logarithm of $\cos \alpha$.
  Since the $\cos$ is equal to the sine of the complementary angle, 
  we find this value in the table in the figure, 
  row 0, column 
  \emph{Differenti\ae}: $1.8427293$ 
  (all values are multiplied by $10^7$).
  The sought value is the exponential of this 
  logarithm, and we can calculate it by finding 
  the same value in the \emph{Logarithmi} column
  and taking the corresponding value in the 
  \emph{Sinus} column. The value is between the values 
  of row 6 and row 7. More precisely, 
  it differs from the value in row 7 
  $\frac{3842}{18143}$ times the difference between the value 
  of row 6 and the value of row 7. 
  Applying the same ratio to the values found 
  in the \emph{Sinus} column, we obtain
  by interpolation that from the value 
    of row 7 (equal to 1.584453), we subtract 608, thus obtaining $\tg \alpha
  \approx 0.1583845$,
  which differs from the exact value 
  by less than $10^{-7}$.
  }
  \label{fig:napier}
\end{figure}

\label{nota:Nepero}%
\index{Napier!John}%
\index{Nepero}%
\label{Euler!Leonhard}%
\label{Eulero}%
Nepero is the Italianization of the name of the
Scottish mathematician \emph{John Napier} (1550-1617),
who, to compile the tables of the sine function
with a precision of 7 decimal digits, was the first to introduce 
the logarithm function (see figure~\ref{fig:napier}).
It is interesting to note that Napier directly defined 
the natural logarithm without introducing the number $e$ 
and without making any connection to the exponential function.
% https://archive.org/details/johnnapierinvent00hobsiala/page/18/mode/2up

\label{note:Bernoulli}%
\index{Bernoulli!Jacob}%
The identification of the constant $e$
is due to \emph{Jacob Bernoulli} (1655--1705) in 1683.
Bernoulli wondered what the effective annual interest
is obtained from a capital that yields a daily income.
If I invest a capital $c$ at an annual interest rate 
of $r$, at the end of the year I receive back 
the capital $c$ plus an interest $rc$.
If I divide the gain $rc$ by the number of days in 
a year, I can consider having had a daily gain
equal to $\frac{r}{365} c$.
If the daily gain $\frac{r}{365} c$ were returned 
immediately every day of the year, I could reinvest the gain immediately, 
increasing the capital during the year.
In this way, if I start with a capital $c$, each day my capital 
would increase by a factor of $1+\frac{r}{365}$, and
at the end of the year, I would have a gain
equal to $\enclose{1+\frac{r}{365}}^{365} \approx e^r$ 
times the initial capital (compound interest).

The name $e$ was introduced by Euler (Leonhard Euler 1707-1783).
%http://eulerarchive.maa.org//docs/originals/E853.pdf

\subsection{The Fundamental Theorem of Algebra}

The formula for determining the roots of a second-degree polynomial 
was already known to the Babylonians, and although complex numbers were not
known
(nor were negative numbers), the same formula applies 
in the general case and provides the complex roots.

The study of equations of degree higher than the second developed 
in the 1500s by Italian mathematicians Scipione del Ferro, 
Gerolamo Cardano, Niccolò Tartaglia, and Lodovico Ferrari.
They determined the formulas for finding the roots 
of third and fourth-degree polynomials. 
When applying the formula for the third-degree equation, 
one might need to calculate the square root of a negative number
even in situations where all solutions ultimately turn out to be real.
It is precisely in this context that complex numbers arise, as 
a mere mathematical artifice used to complete the calculation.

Subsequently, Paolo Ruffini (1765-1822) and Niels Abel (1802-1829) demonstrated
that equations
of degree higher than the fourth do not admit a formula 
expressible by radicals. 
The criterion for determining whether a polynomial
admits solvable formulas is due to the French mathematician Évariste Galois
(1811-1832), who is considered the founder of group theory.

The fundamental theorem of algebra, namely the existence of complex roots 
for polynomials of any degree, was demonstrated by 
the German mathematician Carl Friedrich Gauss (1777--1855). 
This theorem marked a turning point in mathematical thought as 
for the first time, importance was given to an abstract result of existence 
unrelated to a solvable formula. 
The theorem was not at all obvious, considering that both Leibniz 
and Nikolas Bernoulli believed they had found fourth-degree polynomials 
that could not be factored, contrary 
to theorem~\ref{th:fattorizzazione_polinomio_reale}.

\subsection{The Isoperimetric Problem}

\label{note:isoperimetrico}%
\index{problema!isoperimetrico}%
It was probably already known to the Greeks in the 9th century B.C.\ that 
the circumference, among all plane lines of fixed length, 
is the one that encloses the largest area (isoperimetric property).
In the Aeneid, Virgil recounts that Dido faced 
a similar problem in the founding of Carthage.
\index{Didone}%
\index{isoperimetrico!problema}%
\index{problema!isoperimetrico}%
\index{problema!di Didone}%
Ascending to dimension $n=3$, it is natural to imagine that 
the surface enclosing the largest volume with fixed area 
is the sphere. 
This explains why a drop of water, in the absence of gravity,
should assume a spherical shape: due to surface tension, 
the surface of the drop tends to have the minimum area.
Euler (\emph{Leonhard Euler} 1707--1783)
\index{Euler!Leonhard}%
\index{Eulero}%
thought he had
found a proof
based on this result:
given any closed surface in space,
if it is not a sphere, it is possible
to make a small modification to obtain a new surface
that encloses the same volume but has a strictly smaller area.
Dirichlet (1805--1859) noticed a fundamental problem 
in Euler's proof and in subsequent formalizations by Steiner:
Euler demonstrated that the minimum cannot be a surface other than the sphere,
but did not prove that the minimum exists.
In fact, Euler shows that there cannot be a minimum that is not the sphere, 
but does not prove that the minimum exists, so it is not certain that the sphere
is the minimum:
there could be infinite surfaces $S_1, S_2, \dots, S_n, \dots$ 
each with an area smaller than that of the sphere and each with an area greater
than the next:
$A(S_{n+1})< A(S_n)$. 
To conclude Euler's reasoning, it was necessary to prove 
that in such a situation, the surfaces $S_n$ must, in some sense, 
tend to the spherical surface and that their areas must
tend to the area of the sphere. 
Karl Weierstrass (1815--1897)
\index{Weierstrass!Karl}%
undertook to refound analysis to formalize these proofs.
The proof of the isoperimetric property of the circle and the sphere was 
formalized by Hermann Schwarz, 
\index{Schwarz!Hermann}%
a student of Weierstrass's school.
The case of the circle is usually attributed to Steiner (1796--1863) because 
although he, like Euler, did not notice the problem of the existence of the
minimum,
he had produced all the necessary tools to complete the proof. 

The formulation of the isoperimetric problem in its most general form 
(that is, for sets that are not necessarily regular, and in any dimension)
is due to the Italian mathematician Ennio De Giorgi (1928--1996)
\index{De Giorgi!Ennio}%
who formalized the concept of perimeter introduced
by Renato Caccioppoli (1904--1959).%
\index{Caccioppoli!Renato}%
\index{Caccioppoli!perimetro}%